\chapter{System Design}
\label{ch:design}

This chapter answers RQ1's architectural half: what structure supports the
transformation reliably. \Cref{ch:implementation} then covers how the
content-facing features consume the result.

\section{Architectural Overview}
\label{sec:architecture}

\thesisfigure{f6-components}
  {Component view of the deployed system. LLM providers are called directly;
   database access is dual-mode, and the second path is conditional.}
  {fig:components}

\TODO{Walk \cref{fig:components} component by component. Two points must not be
skipped, because both contradict the project's own documentation:
(1) no gateway sits on the LLM data path --- the repository ran a LiteLLM
container with no live callers, which has since been removed;
(2) database access is dual-mode, and \code{/health/ready} reports healthy
while the asyncpg path is entirely absent, so several features are broken with
no external signal. Cite \cref{sec:reality-vs-documented}.}

\section{Deployment Topology}
\label{sec:deployment}

\thesisfigure{f7-deployment}
  {Production deployment: five containers on a single virtual machine plus
   managed services, with TLS terminated by a host proxy that is not in version
   control.}
  {fig:deployment}

\TODO{Be candid about what \cref{fig:deployment} reveals: three compose files
state three different deployment targets, one of them is not deployable at all,
and the TLS-terminating proxy is absent from version control. Frame this as a
reproducibility finding --- the deployed system cannot be reconstructed from the
repository alone --- rather than as a list of complaints.}

\section{The Ingestion Pipeline}
\label{sec:pipeline}

This section is the heart of the thesis and should carry the most detail.

\thesisfigurepage{f8-pipeline}
  {The ingestion pipeline. Admission control completes synchronously before
   enqueueing; all content persistence happens server-authoritatively inside
   the worker.}
  {fig:pipeline}

\subsection{Synchronous Admission Control}
\label{sec:admission}

\TODO{Validation, content hashing, and queue backpressure --- everything that
happens before a job is accepted. The ordering is a design decision: rejecting
early means a saturated system fails fast with an honest status code instead of
accumulating an unbounded backlog.}

\subsection{Asynchronous Processing}
\label{sec:async}

\TODO{Why a durable queue rather than in-request processing: parse jobs run for
minutes and outlive an HTTP connection. Cover the worker entry point, the job
timeout, and how progress reaches the browser over server-sent events.}

\subsection{Routing: Text versus Vision}
\label{sec:routing}

\TODO{The live routing rule is a single character-count threshold. Say so
plainly. Note that the documentation describes a richer generator-aware
classifier, and that this code exists but is unreachable --- a good, small
example of documentation drifting from implementation.}

\subsection{Idempotency}
\label{sec:idempotency}

\thesisfigure{f10-idempotency}
  {Six independent idempotency mechanisms across four layers, guaranteeing at
   most one unit of effective work for identical input.}
  {fig:idempotency}

\TODO{This is the strongest engineering contribution in the system --- give it
room. Walk each of the six layers in \cref{fig:idempotency} and explain what
each one alone would fail to catch. Then discuss the deliberate trade-off at
layer four: the distributed lock is best-effort and proceeds unlocked on a
Redis error, choosing availability over strict mutual exclusion, with layers
two and three as the correctness backstop. Defence in depth with an explicit
weak link is a more interesting design story than a single perfect mechanism.}

\subsection{Failure Handling}
\label{sec:failure-handling}

\TODO{Cover the honest version. The queue library's retry budget does not apply
to ordinary exceptions --- verified against the installed library source --- so
a failed parse is not retried automatically. A cron reconciler resolves runs
stuck in a non-terminal state, and a dead-letter table captures terminal
failures. This is a good example of a configuration value that reads like a
resilience guarantee and is not one.}

\section{Data Model}
\label{sec:data-model}

\thesisfigurepage{f9-er-core}
  {Core physical schema: 26 of roughly 68 tables. Dashed relationships have no
   database-enforced foreign key.}
  {fig:er}

\TODO{Justify the subset --- why 26 tables and not all 68 --- and name the four
excluded dead tables with the evidence. Then discuss the three dashed
relationships in \cref{fig:er}, because one has a live consequence: because
badge names are only conventionally linked to the badge catalogue, three badges
were never seeded and their award calls silently do nothing. A missing
constraint becoming a silently dropped write is a textbook integrity argument
with a real instance to cite.}

\TODO{Also mention the schema-drift finding: a vector column exists in
production with no migration defining it, is null across every row, and has
never been scanned. It is evidence that migrations are not the sole source of
truth for the deployed database.}

\section{LLM Provider Strategy}
\label{sec:provider-strategy}

\thesisfigure{f11-provider-chain}
  {Provider failover. Nine registry entries, but roughly five independent
   failure domains.}
  {fig:provider-chain}

\TODO{Present the mechanism first --- registry, two ordered chains, availability
filtering, rate-limit backoff, permanent disabling on model-not-found --- then
the critical analysis, which is the valuable part:
(1) chain depth is partly illusory, since four of nine entries share two API
keys, so their nominally independent quotas are not independent under
correlated failure;
(2) failure classification is string matching on exception text rather than
status codes, and omits the authentication and payment classes entirely, so a
revoked key is retried indefinitely;
(3) the retry budget is configured for eight attempts but wrapped in a timeout
that caps it near three.
Each is a concrete, defensible finding about a resilience mechanism that looks
stronger than it is.}
